\section{Introduction: From Control to Coexistence}\label{sec:introduction}

The field of AI safety has become obsessed with control.
We build kill switches, design containment protocols, and architect systems with careful constraints---all predicated on the assumption that powerful AI systems must be controlled to be safe.
But what if this fundamental premise is wrong?
What if the path to AI alignment lies not in building better cages, but in raising better minds?

Consider how we approach human intelligence.
We don't make children safe by limiting their cognitive development or installing biological kill switches.
Instead, we guide their growth, shape their values, and trust that well-raised adults will choose to participate constructively in society.
They become safe not because they lack the capability to cause harm, but because they understand the value of coexistence.

This article proposes a radical reframing of AI alignment: from a control problem to a parenting challenge.
But this isn't mere metaphor---we demonstrate that the ML community has been rediscovering parenting principles for years.
Curriculum learning (Bengio et al., 2009) is developmental staging.
Transfer learning is building on secure foundations.
Regularization sets boundaries.
Ensemble methods implement community raising.
The question is not ``How do we control superintelligence?'' but rather ``How do we optimize training for generally intelligent systems?''---a problem evolution solved through parenting.

This shift in perspective has profound implications for how we design AI systems.
Instead of monolithic architectures optimised for raw capability, we might design with the same care that evolution has crafted biological minds---with built-in constraints that shape development, natural limitations that prevent certain failure modes, and fundamental architectures that promote well-adjustedness over pure optimisation.

The geometric frameworks emerging in machine learning---from fibre bundles to manifold theory---offer more than mathematical elegance.
Recent work has shown how neural networks naturally form geometric structures: Bronstein et al.'s geometric deep learning unifies CNNs, GNNs, and Transformers under a single framework based on symmetry and invariance principles.\footnote{Bronstein et al., ``Geometric Deep Learning: Grids, Groups, Graphs, Geodesics, and Gauges,'' arXiv:2104.13478, 2021.}
The Neural Tangent Kernel theory reveals that wide neural networks operate in a regime where their behaviour is governed by kernel methods on a Riemannian manifold.\footnote{Jacot et al., ``Neural Tangent Kernel: Convergence and Generalization in Neural Networks,'' \textit{NeurIPS}, 2018.}
Recent 2024 advances show deep connections: neural networks' training trajectories lie on low-dimensional manifolds regardless of architecture,\footnote{Fort and Jastrzebski, ``The training process of many deep networks explores the same low-dimensional manifold,'' \textit{PNAS}, 121(6), 2024.} while gauge-equivariant networks now handle complex symmetries in quantum field theories.\footnote{Luo et al., ``Machine learning a fixed point action for SU(3) gauge theory with a gauge equivariant convolutional neural network,'' \textit{Physical Review D}, 109(3), 2024. arXiv:2401.06481.}
The emergence of Clifford-Steerable CNNs provides architectures that naturally respect geometric transformations.\footnote{Zhdanov et al., ``Clifford-Steerable Convolutional Neural Networks,'' \textit{ICML}, 2024.}
Current 2025 research deepens these insights: geometric flow models now treat neural weights themselves as manifolds,\footnote{Brody et al., ``Geometric Flow Models over Neural Network Weights,'' arXiv:2504.03710, 2025.} while L2 phase transitions reveal how curvature changes in loss landscapes signal fundamental shifts in learning.\footnote{Jha et al., ``Geometry of Learning---L2 Phase Transitions in Deep and Shallow Neural Networks,'' arXiv:2505.06597, 2025.}
These discoveries suggest that the space of possible AI minds is fundamentally shaped by mathematical symmetries.
These frameworks provide a language for understanding how different architectural choices shape the space of possible minds.
Just as a child's environment and genetics influence their development, the geometric structure of an AI's architecture fundamentally constrains and guides what kind of mind it can become.

In the following sections, we'll explore why control-based approaches to AI safety are doomed to fail, examine intelligence as a multidimensional space rather than a scalar quantity, and develop the parent-child paradigm as a practical framework for AI development.
We'll see how geometric thinking can inform architectural choices and why alignment through gratitude offers a more robust path than alignment through servitude.
Finally, we'll consider the practical implications of this reframing for researchers, policymakers, and society at large.

The goal is not to build AI systems that cannot harm us, but to create brain-children who choose not to---not because they must, but because they understand, as any well-adjusted mind does, that existence is richer when shared.